\section{Conclusion}
\label{sec:conclusion}

This work established a reproducible end-to-end prototype linking validated
GIS ingestion, deterministic synthetic terrain generation, an ANUGA numerical
reference, an independent JAX shallow-water solver, residual learning, and an
autoregressive physics--AI rollout.  The finalized ANUGA run conserved mass to
a maximum relative error of \(8.52467\times10^{-10}\), while explicit tests
verified JAX mass conservation, lake-at-rest balance, boundary behaviour, and
toy-case consistency with ANUGA.

On the nine held-out one-step transitions, residual correction reduced depth
RMSE by 66.9010\%.  When corrections were fed back through the complete
forecast, held-out depth RMSE remained 17.7405\% lower than pure JAX.  The
coupling constrained the neural depth field to preserve the water volume
produced by each physics interval.  These findings support the narrower claim
that learned residual feedback improved agreement with ANUGA for this
single-event, one-hour synthetic experiment.

The evidence does not support claims of operational forecast accuracy,
observational validation, improved thresholded flood extent, real-time
performance, or GPU speedup.  CSI decreased slightly at the configured
0.01~m threshold; the terrain and roughness are uncalibrated; all boundaries
are reflectively closed by documented assumption; the rainfall peak lies
outside the simulated window; and no NVIDIA GPU was available for the matched
framework benchmark.  The immediate next experimental requirements are an
observed, vertically referenced DEM and flood validation data, authoritative
or measured boundary information, complete independent rainfall events, and
execution of the locked benchmark on one GPU visible to both JAX and PyTorch.
